
% ================================
% SFU CMPT 276 PHASE 2 REPORT
% ================================

\documentclass[12pt,letterpaper]{article}

% Packages
\usepackage[margin=0.5in]{geometry}

% Title of tex report
\title{Phase 2 Report}

\author{}

\begin{document}

\maketitle

% Information about Group 2
\begin{tabular}{ll}
\textbf{Project/Course:} & CMPT 276 D300, Phase 2 \\
\textbf{Team Name:}      & Group 2 \\
\textbf{Team Members:}   & Harman Bains, Nathaniel Chen, Jayden Pai, \\
                         & Komal Kaur, Maryna Almeida
\end{tabular}

% Body of report:

\section{Overall Game Plan}
The phase 2 game plan was to begin with creating method declarations inside of each class. By doing this, we were able to utilize abstraction of these classes to create our own methods without worrying about how other classes worked in the background. Afterwards, our plan was to create the logic for and between these classes. As a result, we were able to implement the Java Swing portion of the project with ease, as all of our methods and logic were already there. Furthermore, our plan was to proactively bug-fix while implementing the Java Swing animations, as we were now able to visually see where rewards, entities and barriers were, and were able to see how the logic of these classes changed the panels of our game, and fix it if it was not desirable.

\section{Management Process}
The phase 2 management process was noted to be split into two parts. The first part was figuring out among group members which parts they would like to do in the coding implementation. It was considered to be finalized by the TA meeting day. The second part was executing tasks by each group member a few days prior to the final deadline. This allowed people to work on the phase 2 report which required less time. Leading to how the division of roles and responsibilities was done which is towards conducting an even split. Where, each person in the group decide which parts to code through the UML diagram to select that would result in even work to distribute based on splitting total files needed. The same rules applied with the report also and larger coding structures which allowed to have collaboration among group members for efficiency of the time constraint. 

\section{Biggest Challenges Faced}
One of the first major challenges faced for the group was managing time. It can be difficult to find time to communicate in meetings or complete tasks among five people whether it was due to school, work, and other outside life circumstances. The second challenge was solving group issues which were unexpected. An example includes missing members in meetings for role dispersion where it was resolved for the person to receiving the remaining tasks left by default. A third challenge was even work distribution among each individual in the group. It was the goal to strive for evenly split tasks, but despite conducting this method it may still lead to a few people having a higher workload within their role compared to others that needed to be resolved. The last major challenge was understanding the direction of the project. It is not a simple task to start from the scratch especially for phase two, so managing how to start and develop work was a key to success in order for completion.

\section{Modifications \& Adjustments}
One of the modifications made to the initial design was adding more methods to the classes which were laid out in the UML class diagram. This includes examples of adding the get methods such as get cell type as there was set methods also, but they were not binding data within the class known as encapsulation. Another adjustment created was implementing a chest java file as it was missing from phase 1. It went unnoticed in the UML and required to connect the rewards folder in terms of the coins which needed to be used for the opening of a chest. Moreover, there was a necessity of adding a UI folder into the source folder. The reason for this adjustment was to add the game view java file which involved the spirits process of animation especially the game functionality such as the difference between a opened and closed chest which required gather image png paths. Additionally, the UI folder also contained java files for pages including home, winning, losing, and tutorial which were needed for the start-up of the game and scenarios following it as well. This leads to the last modification which was adding an asset folder. The purpose of it was to hold image pngs for the game and the map structure which was not taken into consideration of executing during phase 1.

\section{External Libraries Utilized}
The first library used which may not necessarily be external is importing java swing. The purpose of it was to include UI components which examples are in the main class that needed to use the J frame in order to have a window to hold the GUI components together in place. Another reason was for the J panel in game view class that required to have the role of organizing the GUI components within the window. A second library used which again may not be applicable as external is importing java-awt known as the abstract window tool-kit. This library is specifically regarding tools for the GUIs which the focus was on java-awt events which are for the states altering in a component such as a button. Examples for its purpose of utilizing it was in game view which the java-awt action event as the name suggests proceeds followed with the action listener taking it in and responding as a result. Overall, there were key libraries used for the GUI which as mentioned are not considered external, but can be applied for as a replacement in this case.

% ================================
\end{document}
% ================================
```
